\documentclass[10pt,a4paper]{amsart}
\usepackage[margin=19mm]{geometry}
\usepackage{amsmath,amssymb,mathtools}
\usepackage[scr=rsfs,cal=euler]{mathalfa}
\usepackage{xfrac}
\usepackage[unicode,hidelinks,bookmarks=false]{hyperref}
\newcommand{\ie}{i.e.\ }
\newcommand{\cf}{cf.\ }
\newcommand{\resp}{respectively\ }
\newcommand{\Subsection}{\subsection}
\providecommand{\nobreakdash}{\nobreak\hbox{-}}
\setlength{\emergencystretch}{2em}
\multlinegap0pt
\usepackage[shortlabels]{enumitem}
\usepackage{xcolor}
\usepackage{amssymb}
\newtheorem{theorem}{Theorem}
\newtheorem{lemma}{Lemma}
\newcommand{\E}{\mathbb{E}}
\title{Optimal Feedback Communication with Information Maximization and Distortion Minimization}
\author{Aolin Xu}
\date{Source: arXiv:2606.09698v1 (CC BY 4.0)}
\begin{document}
\maketitle

\noindent\emph{Source and licence.} Optimal Feedback Communication with Information Maximization and Distortion Minimization. Version arXiv:2606.09698v1 (CC BY 4.0), distributed by arXiv under the Creative Commons Attribution 4.0 International licence, http://creativecommons.org/licenses/by/4.0/, as recorded in the arXiv licence metadata for this exact version and shown on the abstract page https://arxiv.org/abs/2606.09698v1. Full licence notice: \texttt{licenses/arxiv-2606.09698-CC-BY-4.0.txt}.\par\smallskip

\noindent\textit{Editorial scope.} Counted unit: the article's lemma lm:E[W|Y] (source0.tex lines 267--282). The article's complete original proof of it is delivered with it (source0.tex lines 283--329, one \texttt{proof} environment of 47 lines). No mathematics is written, repaired or re-derived: the statement and the proof are the article's own lines, in the article's own notation, and no retained line is substituted or re-flowed.  Retained uncounted as the source-local context the proof needs: lines 205--216; lines 239--241.  Nothing was omitted from any retained span, so the package carries no omission record.  No retained line contains a citation, so the wrapper carries no bibliography and no bibliography entry.  No retained line contains a figure environment, an image-inclusion command or drawing code, so no image is delivered and the package declares no asset dependency.  Editorial formatting only, in addition: an AMS \texttt{amsart} preamble that restates the article's own declarations and macros used by the retained lines, namely the environments \texttt{theorem}, \texttt{lemma} on separate counters, together with the standard packages those lines need.  The retained environments print in this document as Theorem 1, Lemma 1 in order of appearance, whereas the article prints the same items with the numbering of its own full document.  The complete native source of the article as distributed in the arXiv e-print for this exact version is bundled unchanged as \texttt{sources/202422/source0.tex}.
\begin{theorem}\label{th:suff_info_max}
If the sequence of encoding functions $(\varphi_t)_{t=1}^n$ satisfy the following three conditions for all $t=1,\ldots,n$: 
\begin{enumerate}[leftmargin=*]
    \item 
    $\varphi_t$ is a deterministic function, meaning that $X_t$ is uniquely determined by $(W, Y^{t-1})$;
    \item 
    the marginal distribution of $X_t$ is a capacity-achieving distribution $P_X^*$ of the channel;
    \item 
    $X_t$ is statistically independent of $Y^{t-1}$;
\end{enumerate}
then $I(W;Y^n)$ achieves the maximal value $nC$.
\end{theorem}

\begin{theorem}\label{th:nece_info_max}
Under the assumption that the channel $P_{Y|X}$ is input-identifiable, the three conditions listed in Theorem~\ref{th:suff_info_max} are also necessary for the sequence of encoding functions $(\varphi_t)_{t=1}^n$ to satisfy for all $t=1,\ldots,n$ in order to have $I(W;Y^n) = nC$.
\end{theorem}

\begin{lemma}\label{lm:E[W|Y]}
Under the assumption that the channel $P_{Y|X}$ is input-identifiable, with a sequence of encoding functions belonging to $\mathsf\Phi$, we have for each $t=1,\ldots,n$,
\begin{align}
& \E[\E[W|Y^{t}]^2] 
% = \nonumber \\
% &\qquad 
=\sum_{y^{t-1}}P_{Y^{t-1}}(y^{t-1}) \E[\E[W|Y^{t-1}=y^{t-1}, Y_t]^2] \label{eq:lm_E[W|Yt]_1}
\end{align}
and
\begin{align}
& \E[W|Y^{t-1}=y^{t-1}, Y_t = y_t] = 
%\nonumber \\
%&
\sum_{x_t}\frac{P_{X}^*(x_t) P_{Y|X}(y_t|x_t)}{P_{Y}^*(y_t)} \E[W|Y^{t-1}=y^{t-1}, X_t=x_t] . \label{eq:lm_E[W|Yt]_2}
\end{align}
\end{lemma}

\begin{proof}
Under the assumptions that the channel $P_{Y|X}$ is input-identifiable and the encoder belongs to $\mathsf\Phi$, Theorem~\ref{th:nece_info_max} implies that $X_t$ and hence $Y_t$ are statistically independent of $Y^{t-1}$.
Equation~\eqref{eq:lm_E[W|Yt]_1} follows from
\begin{align}
% &\quad 
\E[\E[W|Y^{t}]^2]  
%\nonumber \\
% &
&= \sum_{y^{t-1}}P_{Y^{t-1}}(y^{t-1}) \E[\E[W|Y^{t-1}=y^{t-1}, Y_t]^2|Y^{t-1}=y^{t-1}] \nonumber\\
&= \sum_{y^{t-1}}P_{Y^{t-1}}(y^{t-1}) \E[\E[W|Y^{t-1}=y^{t-1}, Y_t]^2] \label{eq:lm_E[W|Yt]_3}
\end{align}
where \eqref{eq:lm_E[W|Yt]_3} follows from the necessary condition that $Y_t$ is statistically independent of $Y^{t-1}$, as a consequence of the above assumptions.

Equation~\eqref{eq:lm_E[W|Yt]_2} follows from
\begin{align}
% &\quad 
\E[W|Y^{t-1}=y^{t-1}, Y_t=y_t]  
% \nonumber \\
&= \sum_{x_t} P_{X_t|Y^{t-1}, Y_t}(x_t|y^{t-1}, y_t)
% \cdot \nonumber \\
% &\qquad 
\E[W|Y^{t-1}=y^{t-1}, Y_t=y_t, X_t=x_t] \\
&= \sum_{x_t}\frac{P_{X}^*(x_t) P_{Y|X}(y_t|x_t)}{P_{Y}^*(y_t)} \E[W|Y^{t-1}=y^{t-1}, X_t=x_t] 
% \nonumber
\end{align}
where the last step follows from
\begin{align}
% & \quad 
P_{X_t|Y^{t-1}, Y_t}(x_t|y^{t-1}, y_t) 
% \nonumber \\
&= \frac{P_{X_t|Y^{t-1}}(x_t|y^{t-1})P_{Y_t|X_t}(y_t|x_t)}{P_{Y_t|Y^{t-1}}(y_t|y^{t-1})}  \label{eq:lm_E[W|Yt]_5}\\
&= \frac{P_{X_t}(x_t) P_{Y_t|X_t}(y_t|x_t) }{P_{Y_t}(y_t)}  \label{eq:lm_E[W|Yt]_6} \\
&= \frac{P_{X}^*(x_t) P_{Y|X}(y_t|x_t)}{P_{Y}^*(y_t)}  \label{eq:lm_E[W|Yt]_7}
\end{align}
and from
\begin{align}
&\E[W|Y^{t-1}=y^{t-1}, Y_t=y_t, X_t=x_t] = 
% \nonumber \\
% &\qquad 
\E[W|Y^{t-1}=y^{t-1}, X_t=x_t]  \label{eq:lm_E[W|Yt]_8}
\end{align}
where \eqref{eq:lm_E[W|Yt]_5} follows from the Markov chain $Y^{t-1}-X_t-Y_t$; 
\eqref{eq:lm_E[W|Yt]_6} follows from the necessary condition that $X_t$ and $Y_t$ are statistically independent of $Y^{t-1}$ under the assumptions;
\eqref{eq:lm_E[W|Yt]_7} follows from the fact that $P_{Y_t|X_t}=P_{Y|X}$ and the necessary condition $P_{X_t}=P_X^*$ hence $P_{Y_t}=P_Y^*$ under the assumptions; 
and 
\eqref{eq:lm_E[W|Yt]_8} follows from the Markov chain $W - (Y^{t-1},X_t) - Y_t$.
\end{proof}

\end{document}
